% problem-000325 8 镁基储氢晶体
\section*{8. 镁基储氢晶体}
\textit{下列为分问。}

\subsection*{8-2}
将上述氢化物与⾦属镍在⼀定条件下⽤球磨机研磨，可制得化学式为 $\mathrm { M g N i H _ { 4 } }$ 的化合物。X射线衍射分析表明，该化合物的⽴⽅晶胞的⾯⼼和顶点均被镍原⼦占据，所有镁原⼦的配位数都相等。推断镁原⼦在$\mathrm { M g N i H _ { 4 } }$ 晶胞中的位置（写出推理过程）。

\subsection*{8-3}
实验测定，上述 $\mathrm { M g N i H _ { 4 } }$ 晶体的晶胞参数为 646.5 pm，计算该晶体中镁和镍的核间距。已知镁和镍的原⼦半径分别为 159.9 pm 和 124.6 pm。

\subsection*{8-4}
若以材料中氢的密度与液态氢密度之⽐定义储氢材料的储氢能⼒，计算 $\mathrm { M g N i H _ { 4 } }$ 的储氢能⼒（假定氢可全部放出；液氢的密度为0.0708 $\mathrm { g c m ^ { - 3 } } )$
